\subsection{Роль амплитудной и фазовой составляющих комплексного S-параметра}

Как было показано в параграфе 2.2, SFCW-радар на каждой частоте $f_k$ измеряет комплексный коэффициент рассеяния
\begin{equation}
    S_{ij}(f_k) = |S_{ij}(f_k)| e^{i\varphi_{ij}(f_k)}, \qquad i,j \in \{1,2\}.
\end{equation}
Наиболее часто в георадиолокационных системах анализу подвергается коэффициент передачи $S_{21}$, отражающий связь между передающим и приемным портами. Однако возможна работа в конфигурации, когда передача и прием осуществляются с одной и той же антенной, и анализируется коэффициент отражения $S_{11}$. Поэтому далее <<S-параметр>> обозначает обобщенное семейство величин $S_{ij}$.

При конструировании радара важно понимать, какая из компонент комплексного S-параметра --- амплитуда или фаза --- играет ключевую роль в формировании радиолокационного изображения. В современной литературе отмечается, что именно фаза содержит основную информацию о временных задержках, определяющих глубинную структуру среды. Так, в работе \cite{eide2019advanced} продемонстрировано, что даже при полном устранении амплитудных колебаний, путем замены всех $|S(f_k)|$ на единицу, обратное преобразование во временную область сохраняет геометрию B-скана, но с увеличением уровня шумов. С физической точки зрения, такое нормирование эквивалентно предположению, что равная доля излученной энергии отражается обратно на всех частотах, а суммарная энергия во временной области сохраняется. В этом случае сумма квадратов отсчетов в каждом A-скане одинакова:
\begin{equation}
    \sum_n |s_{\mathrm{in}}[n]|^2 = 1.
\end{equation}

Тем не менее, в процессе работы было установлено, что в простых случаях, например при наличии единственного отражателя, информация, содержащаяся в амплитудном спектре, также может быть использована для восстановления пространственной структуры. Для анализа этого эффекта рассмотрим модельную задачу: излучающая и приемная точечные антенны расположены на расстоянии $d$ друг от друга, а обе они находятся на расстоянии $L$ от идеально отражающей плоской поверхности. Излученная волна достигает приемной антенны двумя путями: прямым (длина $d$) и отраженным (длина $2L$). Если пренебречь дифракцией и учитывать возможный фазовый сдвиг при отражении $\varphi_{\mathrm{ref}}$, то поле прямого луча в точке приемной антенны можно представить в виде
\begin{equation}
    E_{\mathrm{direct}}(t) = A \sin(kd - \omega t),
\end{equation}
а отраженного луча --- как
\begin{equation}
    E_{\mathrm{refl}}(t) = B \sin(2kL - \omega t + \varphi_{\mathrm{ref}}),
\end{equation}
где $A$ и $B$ --- амплитуды соответствующих компонент, $k = \omega/c = 2\pi\nu/c$ --- волновое число. Тогда результирующее поле на приемной антенне запишется как
\begin{equation}
    E_{\mathrm{tot}}(t) = E_{\mathrm{direct}}(t) + E_{\mathrm{refl}}(t).
\end{equation}

Средняя по времени мощность принимаемого сигнала пропорциональна квадрату полной амплитуды, усредненному по периоду колебаний:
\begin{equation}
    \langle P(\nu) \rangle \propto \langle E_{\mathrm{tot}}^2(t) \rangle_t.
\end{equation}
Используя стандартные тождества
\begin{equation}
    \sin\alpha \sin\beta = \frac{1}{2}\cos(\alpha - \beta) - \frac{1}{2}\cos(\alpha + \beta),
    \qquad
    \langle \sin^2(\theta) \rangle_t = \frac{1}{2},
\end{equation}
и учитывая, что разность фаз $k(d - 2L) - \varphi_{\mathrm{ref}}$ не зависит от времени, можно записать следующее выражение:
\begin{equation}
    \langle P(\nu) \rangle \propto
    \frac{A^2 + B^2}{2} + AB \cos\left[\frac{2\pi\nu}{c}(d - 2L) - \varphi_{\mathrm{ref}}\right].
\end{equation}

Таким образом, наблюдается гармоническая модуляция мощности в частотной области, обусловленная интерференцией между прямой и отраженной волнами. Расстояние между соседними максимумами определяется условием изменения аргумента косинуса на $2\pi$:
\begin{equation}
    \Delta\nu = \frac{c}{2L - d} \approx \frac{c}{2L},
\end{equation}
если $d \ll 2L$. Следует подчеркнуть, что периодическая структура в спектре с характерным периодом $\Delta\nu$ при переходе во временную область дает максимум на задержке
\begin{equation}
    t = \frac{1}{\Delta\nu} = \frac{2L}{c},
\end{equation}
что полностью соответствует ранее сделанным выводам о пиках в A-сканах и глубине отражающих объектов (см. параграф 2.2). Таким образом, даже при отсутствии фазовой информации, по структуре амплитудного спектра можно выявить наличие отражающих границ в пределах длины когерентности сигнала. В диапазоне около 1 ГГц при спектральной ширине каждого тона порядка 1 МГц длина когерентности исчисляется сотнями метров, что позволяет четко наблюдать соответствующие интерференционные осцилляции в лабораторных условиях.

Однако при наличии нескольких границ отражения характер интерференции и структура спектра становятся существенно более сложными и могут приводить к появлению ложных отражений во временной области. В этом случае фаза S-параметра остается основным носителем информации о задержках и глубинах объектов. Тем не менее при качественной аппаратной реализации и достаточном соотношении сигнал/шум информация, содержащаяся только в амплитуде, может быть полезна для выделения простейших объектов, например щелей, расположенных в пределах когерентной длины.

Таким образом, фаза комплексного S-параметра играет первостепенную роль в построении радиолокационного изображения, хотя в ряде приложений допустимо использование одной лишь амплитуды. Наилучшее качество достигается при совместном измерении модуля и фазы. Практическая проверка значимости каждой из компонент проведена в главе 3 на основе лабораторных данных.
